\section{ О распространении упругих волн в напряжённых нелинейных анизотропных средах}

В статье рассматривается влияние начального напряжённого состояния и нелинейных свойств геологических сред на распространение упругих волн. Теория включает:

\begin{itemize}
    \item нелинейные связи между напряжениями и деформациями;
    \item учёт трёхосных напряжений;
    \item анизотропию и неоднородности;
    \item вывод точных уравнений для P, SV и SH волн.
\end{itemize}

Работа развивает классическую теорию упругости, расширяя её возможностью описывать реальные, напряжённые и нелинейные геосреды.

\subsection*{Уравнения движения в анизотропной термоупругой среде}

Рассматривается линейная динамическая задача в анизотропном теле. Для поля смещений $u_i$ уравнения движения имеют вид:
\[
\rho \frac{\partial^2 u_i}{\partial t^2}
=
\frac{\partial \sigma_{ij}}{\partial x_j},
\]
где $\sigma_{ij}$ — тензор напряжений.

Для плоской деформации (P–SV волны) в плоскости $(x_1, x_3)$ система принимает вид:
\begin{align}
\rho \frac{\partial^2 u_1}{\partial t^2} &=
\frac{\partial}{\partial x_1}(\omega_{1111} u_{1,1} + \omega_{1133} u_{3,1}) +
\frac{\partial}{\partial x_3}(\omega_{3111} u_{1,3} + \omega_{3133} u_{3,3}), \\
\rho \frac{\partial^2 u_3}{\partial t^2} &=
\frac{\partial}{\partial x_1}(\omega_{3111} u_{1,1} + \omega_{3133} u_{3,1}) +
\frac{\partial}{\partial x_3}(\omega_{3333} u_{3,3} + \omega_{3311} u_{1,3}),
\end{align}
где $\omega_{ijkl}$ — эффективные упругие коэффициенты, изменённые из-за напряжения и нелинейности.

Для SH волн:
\[
\rho \frac{\partial^2 u_2}{\partial t^2}
=
\omega_{1212} \frac{\partial^2 u_2}{\partial x_1^2}
+
\omega_{3232} \frac{\partial^2 u_2}{\partial x_3^2}.
\]

\subsection*{Представление решения}

Плоские гармонические волны принимают вид:
\[
u_\alpha = A_\alpha e^{i(\omega t - \tau x_1 - v x_3)},
\]
где
\[
\tau = v \sin\theta, \qquad v = \frac{1}{V}, \qquad \theta — угол падения волны.
\]

\subsection*{Фазовые скорости упругих волн}

Подстановка решения в уравнения движения приводит к характеристическому уравнению 4-го порядка:
\[
\begin{aligned}
\big(\omega_{3333}\cos^2 \theta + \omega_{1313}\sin^2 \theta - \rho V^2\big)
\big(\omega_{1111}\sin^2 \theta + \omega_{3131}\cos^2 \theta - \rho V^2\big) \\
-
(\omega_{1133}+\omega_{3113})^2 \sin^2 \theta \cos^2 \theta = 0.
\end{aligned}
\]

Решение даёт скорости продольной (P) и поперечной (SV) волн:
\[
V_{P,SV}^2
=
\frac{\omega_{1111}\sin^2\theta + \omega_{3333}\cos^2\theta \pm Q}{2\rho},
\]
где
\[
Q = \sqrt{
\left(\omega_{1111}\sin^2\theta - \omega_{3333}\cos^2\theta\right)^2
+4(\omega_{1133}+\omega_{3113})^2 \sin^2\theta\cos^2\theta}.
\]

\subsection*{Скорость SH волн}

\[
V_{SH}^2 =
\frac{\omega_{1212} \sin^2\theta + \omega_{3232} \cos^2\theta}{\rho}.
\]

\subsection*{Граничные условия на поверхности раздела}

На границе двух сред должны выполняться:
\begin{itemize}
    \item непрерывность смещений:
    \[
    u_1^{(1)} = u_1^{(2)}, \qquad u_3^{(1)} = u_3^{(2)},
    \]
    \item непрерывность напряжений:
    \[
    \sigma_{13}^{(1)} = \sigma_{13}^{(2)}, \qquad \sigma_{33}^{(1)} = \sigma_{33}^{(2)}.
    \]
\end{itemize}

Эти условия приводят к системе Цептрицца для коэффициентов отражения и преломления.

\subsection*{Нелинейный упругий потенциал}

Используется упругий потенциал 3-го порядка:
\[
\Phi =
\frac12 C_{ijmn}\varepsilon_{ij}\varepsilon_{mn}
+
\frac{1}{6}C_{ijkm}^{(3)} \varepsilon_{ij}\varepsilon_{km}\varepsilon_{mn},
\]
что позволяет учитывать влияние больших напряжений и слабой нелинейности.

В простом анизотропном случае:
\[
\Phi =
\frac12 C_{ijmn}\varepsilon_{ij}\varepsilon_{mn}
+
\frac{a}{3}A_1 + \frac{b}{3}A_2 + \frac{c}{3}A_3,
\]
где $A_1, A_2, A_3$ — инварианты тензора деформации.

\subsection*{Слабо анизотропные напряжённые среды}

Формулы для скоростей выглядят как:
\[
V_P(\theta) = V_{P0}
\left[1 + f_1(\sigma_{ij}) + f_2(\theta, \epsilon, \delta)\right],
\]
\[
V_{SV}(\theta) = V_{S0}
\left[1 + g_1(\sigma_{ij}) + g_2(\theta, \epsilon, \delta)\right],
\]
\[
V_{SH}(\theta) = V_{S0} \sqrt{1 + \gamma \sin^2\theta + h(\sigma_{ij})}.
\]

\subsection*{Численные примеры}

Проведён расчёт скоростей и коэффициентов отражения для различных горных пород:
\begin{itemize}
    \item глина Колорадо,
    \item песчаник Береа,
    \item песчаник Массилон,
    \item белый гранит.
\end{itemize}

Учитывались пять различных трёхосных напряжённых состояний.

Результаты показывают:
\begin{itemize}
    \item изменение амплитуд отражённых волн до 30--40\%;
    \item сильное влияние уровня напряжений на коэффициенты отражения;
    \item наличие противоположных тенденций в зависимости от направления и типа напряжения.
\end{itemize}

\subsection*{Заключение}

Получена трёхмерная теория распространения P, SV и SH волн в напряжённых нелинейных анизотропных средах. Показано существенное влияние начальных напряжений на скорости и амплитуды волн. Разработанные формулы могут быть использованы при интерпретации сейсморазведочных данных и моделировании сложных геологических структур.


\section*{ К теории термоупругих волн}

В работе рассматривается совместное решение линеаризованных уравнений динамической теории упругости и нестационарного уравнения теплопроводности. Показано, что учёт взаимодействия между полем смещений $\mathbf{u}$ и градиентом температуры $\nabla T$ приводит к двум физическим эффектам:
\begin{itemize}
    \item зависимость продольной скорости звука от температуры,
    \item нелинейное (в частности квадратичное по температуре) затухание термоупругих волн.
\end{itemize}

\subsection*{1. Базовые уравнения}

Линеаризованные уравнения для смещений и температуры имеют вид
\begin{align}
    c_{st}^2 \Delta \mathbf{u}' + (1 - 2\sigma)^{-1}(1+\sigma)\nabla(\nabla\!\cdot\mathbf{u}') 
    &= \alpha \nabla T', \\
    \frac{\partial T'}{\partial t} + T \, \nabla\!\cdot\mathbf{u}' &= \chi \Delta T',
\end{align}
где $c_{st}$ — скорость поперечных волн, $\sigma$ — коэффициент Пуассона,  
$\chi$ — температуропроводность, $\alpha$ — коэффициент объёмного расширения.

Решения ищутся в виде волновых зависимостей
\[
\mathbf{u}' = \mathbf{u}_0 e^{i(\mathbf{k}\cdot\mathbf{r} - \omega t)}, 
\qquad
T' = T_0 e^{i(\mathbf{k}\cdot\mathbf{r} - \omega t)} .
\]

\subsection*{2. Продольная скорость звука}

После подстановки плоских волн получено дисперсионное уравнение, приводящее к выражению для частоты:
\[
\omega_0 = k c_{sl}(T),
\]
где продольная скорость упругих волн зависит от температуры как
\[
c_{sl}(T) = c_{sl0}\left(1 - \frac{T}{T^*}\right),
\]
и
\[
c_{sl0} = \sqrt{\frac{E(1-\sigma)}{\rho(1+\sigma)(1-2\sigma)}}.
\]
Таким образом, скорость звука уменьшается линейно с ростом температуры.

\subsection*{3. Затухание термоупругих волн}

Искомое комплексное значение частоты представляется как
\[
\omega = \omega_0 - i\gamma,
\]
где $\gamma(k)$ — коэффициент затухания.  
Его определяет кубическое уравнение
\[
x^3 - A x^2 + B x - a = 0,
\quad 
x = \frac{\gamma}{\chi k^2},
\]
где параметры $A$ и $B$ выражаются через безразмерные величины
\[
a = \frac{T}{T^*}, 
\qquad
b = \frac{\chi k^2}{\omega_0}.
\]

\subsection*{4. Предельные случаи затухания}

\paragraph{1) Длинные волны ($k \ll k^*$).}
\[
\gamma = \chi k^2 \frac{T}{T^*}.
\]

\paragraph{2) Короткие волны ($k \gg k^*$).}
\[
\gamma = \frac{c_{sl}^2 T}{\chi T^*} = \text{const}.
\]

Граница между режимами определяется
\[
k^* = \frac{c_{sl}}{\sqrt{\chi}}.
\]

\subsection*{5. Температурная зависимость затухания}

Так как $\alpha(T)\propto T$ для металлов, то окончательно
\[
\gamma(T) \propto T^2.
\]

\subsection*{Выводы}

\begin{enumerate}
    \item Скорость продольного звука уменьшается с ростом температуры:  
    \[
    c_{sl}(T) = c_{sl0}\left(1 - \frac{T}{T^*}\right).
    \]
    \item Затухание термоупругих волн растёт квадратично по температуре:
    \[
    \gamma \sim T^2.
    \]
    \item При больших $k$ затухание выходит на постоянное значение (режим насыщения).
\end{enumerate}



\section*{Физические и геологические основы сейсморазведки}

\subsection*{Основы теории упругости}  
Геологические среды приближённо рассматриваются как упругие, поэтому теория распространения сейсмических волн базируется на классической теории упругости. При деформации частицы среды смещаются, что характеризуется вектором смещения \(U(x,y,z)\). Деформацию описывают через нормальные и сдвиговые компоненты, всего шесть компонент. Напряжения (силы на единицу площади) вызывают обратимые (упругие) деформации в упругих телах.

\subsection*{Упругие волны в безграничных средах}  
Рассматривается распространение волн в бесконечных (безграничных) однородных средах. Здесь важны волны продольного и поперечного типов, их скорость определяется упругими свойствами среды (модулями упругости).

\subsection*{Упругие волны в слоистых средах}  
В реальной геологической среде часто встречаются слоистые структуры. При этом волны испытывают отражение, преломление и рассеяние на границах слоёв, что существенно влияет на их распространение. Тип распространения волн зависит от контраста свойств слоёв (плотность, модули упругости и др.).

\subsection*{Особенности распространения сейсмических волн в реальных средах}  
\begin{itemize}  
  \item \textbf{Скорости сейсмических волн.} В реальных породах скорости волн зависят от плотности и упругих свойств конкретного геологического материала.  
  \item \textbf{Поглощение (затухание) волн.} В реальных средах волны теряют энергию по мере распространения. Коэффициент поглощения зависит от частоты сигнала и физических свойств породы (пористость, трещиноватость, водонасыщенность и др.).  
\end{itemize}

\subsection*{Выводы}  
\begin{itemize}
    \item Теория упругости — фундамент для описания сейсмических волн в геологии.
    \item В идеализированных (однородных) средах волны распространяются просто, но в природных условиях (слоистость, неоднородности) возникают сложные эффекты.
    \item Атенюация волн — важный фактор в сейсморазведке, так как ограничивает дальность эффективного распространения и влияет на точность интерпретации данных.
\end{itemize}


\section*{Термоупругие волны и скорость их распространения в динамической задаче взаимосвязанной термоупругости}

\author{А.Д. Шамровский, Г.В. Меркотан}

\subsection*{Аннотация}
Исследуется влияние температуры на скорости распространения тепловых и механических волн в полупространстве в рамках обобщённой динамической задачи термоупругости. Взаимосвязанное поведение температурного и механического полей приводит к изменению характеристик распространения волн по сравнению с классической теорией упругости.

\subsection*{Математическая модель}
Для описания термоупругих процессов используется система уравнений:

\begin{align}
\rho \frac{\partial^2 u}{\partial t^2} &= \frac{\partial \sigma}{\partial x}, \\
\sigma &= (\lambda + 2\mu) \frac{\partial u}{\partial x} - k \rho S, \\
S &= \eta^{-1}(T - T_0) + k \eta^{-1} \frac{\partial u}{\partial x}, \\
-q &= \rho T_0 S, \\
\tau \frac{\partial q}{\partial t} + q &= -\kappa \frac{\partial T}{\partial x}.
\end{align}

После преобразований система приводится к волновым уравнениям:

\begin{align}
\frac{\partial^2 \sigma}{\partial t^2} - a_1^2 \frac{\partial^2 \sigma}{\partial x^2} &= d_1 \frac{\partial \sigma}{\partial x} + d_2 \frac{\partial T}{\partial x}, \\
\frac{\partial^2 T}{\partial t^2} - a_2^2 \frac{\partial^2 T}{\partial x^2} &= -d_1 \frac{\partial \sigma}{\partial x} - d_2 \frac{\partial T}{\partial x}.
\end{align}

Где $a_1$, $a_2$ — термоупругие скорости распространения волн, отличающиеся от классических упругих скоростей $c_1$, $c_q$:

\begin{align}
c_1 &= \sqrt{\frac{\lambda + 2\mu}{\rho}}, \quad
c_q = \sqrt{\frac{\mu}{\rho}}.
\end{align}

\subsection*{Результаты}
Проведён численный анализ для различных материалов. Получены значения $a_1$ и $a_2$, учитывающие тепловое влияние. Например, для алюминиевого сплава:

\begin{align*}
c_1 &= 6.1802 \cdot 10^3 \, \text{м/с}, \quad
c_q = 4.4358 \cdot 10^3 \, \text{м/с}, \\
a_1 &= 1.03381, \quad
a_2 = 0.70128 \quad \text{(в безразмерных единицах)}.
\end{align*}

\subsection*{Заключение}
Разработанная модель позволяет учитывать взаимное влияние тепловых и механических волн. Полученные термоупругие скорости важны для анализа поведения материалов при динамических нагрузках и тепловом воздействии.


\section*{Компьютерное моделирование сейсморазведки с трещиноватыми включениями}

\textbf{Авторы:} В.И. Голубев, Н.И. Хохлов, И.Б. Петров (МФТИ)

\subsection*{Аннотация}
Рассматривается задача численного моделирования распространения сейсмических волн в геологических средах с явно заданными трещиноватыми включениями. Предложен сеточно-характеристический метод решения гиперболической системы уравнений упругости на гексаэдральных сетках. Метод позволяет точно учитывать границы геологических слоёв и трещин, обеспечивая высокую точность моделирования сейсмического отклика.

\subsection*{Математическая модель}
Используется модель линейно-упругого изотропного тела. Основные уравнения:

\begin{equation}
\rho \frac{\partial V_i}{\partial t} = \frac{\partial \sigma_{ij}}{\partial x_j}, \quad
\frac{\partial \sigma_{ij}}{\partial t} = q_{ijkl} \frac{\partial V_k}{\partial x_l}
\end{equation}

где $\rho$ — плотность, $V_i$ — компоненты скорости, $\sigma_{ij}$ — тензор напряжений, $q_{ijkl}$ — тензор упругих свойств.

Для одномерного уравнения переноса:

\begin{equation}
\frac{\partial u}{\partial t} + a \frac{\partial u}{\partial x} = 0, \quad a > 0
\end{equation}

применяется схема третьего порядка точности:

\begin{equation}
u^{n+1}_m = u^n_m + \frac{\Delta t}{2h} (\delta_1 + \delta_2) + \frac{\Delta t}{6h} (\delta_1 - 2\delta_0 + \delta_2)
\end{equation}

\subsection*{Особенности реализации}
\begin{itemize}
  \item Явное задание трещин реализуется через корректировку узлов сетки.
  \item Используются контактные условия с полным слипанием:
  \[
  v^a = v^b, \quad f^a = -f^b
  \]
  \item Для трещин применяется модель бесконечно тонкой флюидонасыщенной прослойки.
\end{itemize}

\subsection*{Результаты моделирования}
Проведено моделирование участка геологического массива с мегатрещиной. Использован импульс Рикера с частотой 30 Гц. Получены синтетические сейсмограммы и распределения волнового поля. Показано, что метод с явным заданием границ обеспечивает:
\begin{itemize}
  \item более точное моделирование амплитуд (до 30\% разницы по сравнению со сквозным счётом),
  \item возможность идентификации типов волн (продольные, поперечные, Рэлея),
  \item высокую точность при моделировании сложных геологических структур.
\end{itemize}

\subsection*{Заключение}
Метод позволяет эффективно моделировать сейсмическое поле в трещиноватых средах, улучшая интерпретацию данных сейсморазведки. Результаты применимы для планирования полевых работ и повышения точности инверсии.


